% 分类目录：dl
\documentclass[12pt, a4paper]{article}
\usepackage[margin=2.5cm]{geometry}
\usepackage{ctex}
\usepackage{amsmath, amssymb}
\usepackage{listings}
\usepackage{xcolor}
\usepackage{tabularx}

\lstset{
    language=Python,
    basicstyle=\ttfamily\small,
    keywordstyle=\color{blue},
    commentstyle=\color{green!50!black},
    stringstyle=\color{red},
    showstringspaces=false,
    breaklines=true,
    numbers=left,
    numberstyle=\tiny\color{gray},
    frame=single
}

\title{知识点：生成对抗网络 (GAN)}
\author{}
\date{}

\begin{document}
\maketitle

\section{核心定义}
\textbf{中文定义}：生成对抗网络（Generative Adversarial Network, GAN）是一种通过对抗训练学习数据分布的深度学习架构。它由生成器（Generator）和判别器（Discriminator）组成，两者在相互博弈中共同进化：生成器学习制造逼真数据，判别器学习分辨真伪。

\textbf{英文定义}：Generative Adversarial Networks (GANs) are a framework for training generative models based on a game-theoretic scenario. It consists of two networks: a Generator that learns to produce realistic data and a Discriminator that learns to distinguish between real and generated samples.

\textbf{提出时间与作者}：2014年，由 Ian Goodfellow 等人提出。

\textbf{原始关键论文}：\textit{Generative Adversarial Nets} (NIPS 2014)

\section{核心原理：零和博弈}
GAN 的核心在于其独特的目标函数，定义为一个\textbf{极小极大博弈 (Minimax Game)}。

\subsection{目标函数}
$$ \min_G \max_D V(D, G) = \mathbb{E}_{x \sim p_{data}(x)}[\log D(x)] + \mathbb{E}_{z \sim p_z(z)}[\log(1 - D(G(z)))] $$
\textbf{符号说明}：
\begin{itemize}
    \item $D(x)$：判别器，输出 $x$ 来自真实数据的概率。
    \item $G(z)$：生成器，将随机噪声 $z$ 映射到数据空间。
    \item $\max_D$：判别器希望最大化准确率（真实样本输出近 1，生成样本输出近 0）。
    \item $\min_G$：生成器希望最小化判别器的判断力（使得生成样本让判别器输出近 1）。
\end{itemize}

\subsection{训练过程}
训练是一个交替进行的迭代过程：
\begin{enumerate}
    \item \textbf{固定 G，训练 D}：使判别器更强。
    \item \textbf{固定 D，训练 G}：使生成器造出的图更像真的，从而“欺骗”判别器。
    \item \textbf{纳什均衡}：理论上当 $p_g = p_{data}$ 时，判别器输出恒为 0.5，此时达到平衡。
\end{enumerate}

\section{模型结构示意图}
\begin{verbatim}
[ 随机噪声 z ] --> [ 生成器 G ] --> [ 生成样本 G(z) ] 
                                          |
                                          v
      [ 真实样本 x ] ---------------> [ 判别器 D ] --> [ 分类: 真/假 ]
\end{verbatim}

\section{关键代码片段 (PyTorch实现)}
以下为基础 GAN 的核心训练逻辑简化版：
\begin{lstlisting}
import torch
import torch.nn as nn

# 定义损失函数
criterion = nn.BCELoss()

# 1. 训练判别器 D
# 真实数据 Loss
outputs = D(real_images)
d_loss_real = criterion(outputs, torch.ones_like(outputs))
# 生成数据 Loss
z = torch.randn(batch_size, latent_dim)
fake_images = G(z)
outputs = D(fake_images.detach()) # detach防止更新G
d_loss_fake = criterion(outputs, torch.zeros_like(outputs))
# 反向传播 D
d_loss = d_loss_real + d_loss_fake
d_optimizer.step()

# 2. 训练生成器 G
z = torch.randn(batch_size, latent_dim)
fake_images = G(z)
outputs = D(fake_images) # 重新通过最新的D
# 目标：让 D 认为 fake 是 1 (真)
g_loss = criterion(outputs, torch.ones_like(outputs))
# 反向传播 G
g_optimizer.step()
\end{lstlisting}

\section{深度面试题与解答}
\subsection{基础概念}
\textbf{问：为什么 GAN 的对抗博弈比普通网络更难训练？}\\
答：因为 GAN 不是寻找损失函数的全局最小值，而是寻找两者的纳什均衡点点。如果 D 太强，G 的梯度会消失；如果 D 太弱，G 无法学到任何有用的分布。

\subsection{进阶原理}
\textbf{问：什么是模式崩溃 (Mode Collapse)？}\\
答：指生成器由于发现判别器的某个漏洞，倾向于生成极少数甚至单一的、能够成功欺骗判别器的样本，而失去了数据分布的多样性（例如在 MNIST 上只生成数字 1）。

\subsection{工程实战}
\textbf{问：WGAN (Wasserstein GAN) 解决了什么核心问题？}\\
答：解决了原生 GAN 使用 JS 散度在分布不重叠时梯度消失的问题。WGAN 使用 \textbf{EM距离 (Earth-Mover's distance)} 配合 \textbf{Lipschitz 连续性约束 (Weight Clipping)}，提供了更平滑的梯度映射。

\subsection{坑点分析}
\textbf{问：在评估 GAN 生成质量时，为什么不用训练集 Loss？}\\
答：GAN 的 Loss 只能反映 D 与 G 的博弈进度，不能代表图像质量。常用 \textbf{FID (Fréchet Inception Distance)} 或 \textbf{Inception Score (IS)} 来定量评估图像的逼真度和多样性。

\section{GAN 变体对比表}

\begin{table}[h]
\centering
\begin{tabularx}{\textwidth}{|l|X|}
\hline
\textbf{变体} & \textbf{核心改进} \\
\hline
\textbf{DCGAN} & 引入卷积神经网络（CNN）与 Batch Normalization，确立了工程基础。 \\
\hline
\textbf{WGAN / WGAN-GP} & 引入 Wasserstein 距离与梯度惩罚，极大提升了训练稳定性。 \\
\hline
\textbf{CycleGAN} & 实现了无需配对数据的跨域图像风格转换。 \\
\hline
\textbf{StyleGAN} & 引入解耦控制空间，实现了极高分辨率、可控细节的人脸生成。 \\
\hline
\end{tabularx}
\end{table}

\section{参考文献与拓展}
\begin{enumerate}
    \item \textbf{奠基论文}: Goodfellow, I., et al. (2014). \textit{Generative Adversarial Nets}. NIPS.
    \item \textbf{工程标杆}: Radford, A., et al. (2015). \textit{Unsupervised Representation Learning with DCGAN}.
    \item \textit{GAN Lab}: Google 提供的交互式 GAN 训练可视化工具。
\end{enumerate}

\end{document}
